% CLAIM: two placements of the sum, the byte counts fixed before measurement, and the one
% inequality (docs/11-cost-model.md) that decides between them.
\section{Where to compute the sum}
\label{sec:placement}

Each of $D$ devices holds the full weights and evaluates $N/D$ members. After
evaluation every device has the fitnesses of its own members, and the sum of
Eq.~\ref{eq:update} has to be formed somewhere. There are two natural places.

\begin{itemize}
  \item \textbf{Replicated (A).} Every device gathers all $N$ fitnesses, regenerates all
    $N$ perturbations, and computes the whole sum itself. Only scalars move: $8N$ bytes,
    two gathers of $4N$ (the raw fitnesses, and the shaped weights, which the
    implementation gathers a second time). This is the arrangement behind ``ES only
    communicates scalars'' \citep{salimans2017}.
  \item \textbf{All-reduce (B).} Every device sums its own $N/D$ members, and the partial
    sums are added with an all-reduce of a model-sized buffer: $4P+4N$ bytes for $P$
    float32 parameters. Each device does $1/D$ of the sum.
\end{itemize}

Evaluation and the parameter update are the same in both, so the difference in time per
generation is the sum and the collectives:
\begin{equation}
  t_A - t_B \;=\; \Bigl(1-\frac{1}{D}\Bigr)\,C \;+\; T_{\mathrm{ag}}(4N) \;-\; T_{\mathrm{ar}}(4P),
  \label{eq:crossover}
\end{equation}
where $C$ is the time to form the whole sum on one device and $T_{\mathrm{ag}}$,
$T_{\mathrm{ar}}$ are the times of an all-gather and an all-reduce of the given sizes.
All-reduce wins when the right-hand side is positive: when the share of the sum it saves
takes longer than the model-sized all-reduce. The gather is microseconds
(Table~\ref{tab:tb8}), so in practice
\[
  \text{all-reduce (B) wins} \iff \Bigl(1-\frac{1}{D}\Bigr)\,C > T_{\mathrm{ar}}(4P).
\]
$C$ is large when every member carries full-rank noise that has to be read or
regenerated, and small for low-rank noise. $T_{\mathrm{ar}}$ grows with the model and
depends on the interconnect. So the answer should differ by perturbation, by model
size, by device count and by fabric, and there is no single winner to predict.

\paragraph{What was fixed before measuring.} The byte counts and the direction were
committed before the single-node runs: B should win where $C$ is large (dense and seed
noise) and A where it is small (low-rank). The byte model then needed one correction,
from $4N$ to $8N$ for A (the second gather) and from $4P$ to $4P+4N$ for B, measured from
the compiled program at every cell; it changes nothing, since the comparison is still
kilobytes against 100\,MB. Eq.~\ref{eq:crossover} was written after the single-node
runs, with both terms measured separately (Table~\ref{tab:tb8} for the collectives,
Appendix~\ref{sec:timemodel} for $C$), and used once as a prediction, for the host
boundary of Section~\ref{subsec:boundary}.

\paragraph{The endpoints.} A moves the fewest bytes any placement can and repeats the
whole sum; B splits the sum as far as $D$ devices allow and moves a model-sized buffer.
Other placements trade between the two: a reduce-scatter and all-gather with a sharded
update, as in ZeRO \citep{rajbhandari2020zero}, halves B's bytes; compressed collectives
shrink them further; a hierarchical scheme is A within a host and B across hosts. We
measure the two endpoints.

\paragraph{Overlap.} Hiding B's all-reduce behind the next generation's evaluation is not
free: that evaluation needs the updated weights, so it would run at weights one step
stale, which is a different algorithm. The overlap that keeps the algorithm hides the
all-reduce behind B's own local sum, as data-parallel training does behind the backward
pass, and at best replaces $C/D + T_{\mathrm{ar}}$ by $\max(C/D, T_{\mathrm{ar}})$. With
the measured terms, that still leaves A ahead by 0.40 to 1.98\,ms at every low-rank cell
where A wins at width 2048 (0.77 to 2.29\,ms without overlap). This is a bound from
measured terms; we did not build an overlapped implementation.
